\documentclass[11pt]{article}
\usepackage[margin=2cm]{geometry}
\usepackage{graphicx,amsmath,amssymb,booktabs,longtable,hyperref}
\title{axionFEM validation against arXiv:2211.00008\\[2pt]\large ``Detecting axion dark matter beyond the magnetoquasistatic approximation'' (Benabou, Foster, Kahn, Safdi, Salemi 2022)}
\author{axionFEM \texttt{benchmarks/beyond\_mqs\_2211} (generated by \texttt{make\_note.py})}
\date{\today}
\begin{document}
\maketitle

\section{Purpose}
The 2022 paper computed the beyond-MQS response of a toroidal and a solenoidal lumped-element detector with COMSOL
(RF Module): the source impedance $Z_S(f)$ and source voltage $V(f)$ at a lumped readout port, for three boundary
scenarios (unshielded with absorbing boundaries, shielded by a closed superconducting can, shielded with lossy
dielectric inside), from which it derived quality factors $Q_B$, sensitivity ratios $g/g^{\rm MQS}$ and scan
projections. This note recomputes the paper's figures with axionFEM; the Status section lists which are
reproduced, which are not yet computed and which are not reproduced. It is the validation test of axionFEM's full-wave
(beyond-MQS) solver against an independent code on a published problem.

\section{Method}
\paragraph{Paper curves, extracted exactly.} The paper's figures are vector PDFs. \texttt{extract\_paper\_curves.py}
reads each plotted path, calibrates every axis from its tick labels (each label matched to its tick mark; linear or
logarithmic from the label values and the minor-tick spacing) and maps the vertices to data coordinates. The tick
fit residual is reported below each figure; it is $\lesssim 10^{-3}$\,pt, i.e.\ the extraction is exact to the
precision the PDF stores. Legend labels are attached by matching line styles to the legend samples. Pattern-filled
bands (the inductive-tuning regions) are located on a 600\,dpi rasterization (edges to $\pm0.1$\,pt). Independent
check: the projection plots carry $m_a$ (bottom) and $f$ (top) axes calibrated separately; they agree with
$m_a = 2\pi\hbar f$ ($25$\,neV $\leftrightarrow 6.045$\,MHz).

\paragraph{axionFEM reproduction.} Each figure is recomputed from axionFEM's $Z_S(f)$ and open-circuit
$V(f)=-i\omega\Psi(\omega)$ (reciprocity) on the paper's geometry, then passed through the same circuit and scan
formalism as the paper. The overlays show the extracted paper curves in grey and axionFEM in maroon.

\section{The models (from the paper text)}
\begin{itemize}
\item \textbf{Toroid:} $r_{\rm inner}=0.5$\,m, $r_{\rm outer}=1$\,m, $h=0.5$\,m, PEC walls 2.5\,cm thick, a
5$^\circ$ angular gap bridged by a PEC surface containing a small lumped port. $B_0=10$\,T azimuthal, constant
magnitude, filling the toroid and the gap (axionFEM: \texttt{fullwave.source} kind \texttt{toroidal\_box}).
\item \textbf{Solenoid:} $r_{\rm outer}=2r_{\rm inner}=1$\,m, $h=1.5$\,m, the top gapped by a 1\,cm annulus with the
readout bridge and port; $B_0=10$\,T axial, uniform in the bore, zero outside.
\item \textbf{Scenarios:} unshielded (PML in COMSOL; axionFEM: exact radiation exterior); shielded by a closed
superconducting cylinder ($R=1.5$\,m, $H=1$\,m for the toroid; $H=2.5$\,m for the solenoid); shielded with
TACHYON 100G laminate, 10\% (S3: 20\%) of the field volume (toroid), or lining the solenoid exterior.
\item \textbf{Readout and statistics:} resonant LC readout tuned at each frequency, intrinsic resistance such that
$Q=10^6$ in the MQS limit, $T_{\rm LC}=20$\,mK, $\eta_A = 20$ (main text) or 1 (SM), optimally coupled amplifier;
one-year scan from $f_{\rm MQS}=150$\,MHz downward following DFSZ; $Q_B=\omega/\delta\omega$ from the FWHM of the
response; broadband readout with $S_{VV}^{\rm BA}=\omega^2L_P$, $S_{II}^{\rm imp}=1/L_P$.
\end{itemize}

\paragraph{Not specified in the paper (the reproductions assume values for these; see the status table).}
TACHYON 100G $\varepsilon_r$ and $\tan\delta$ and the placement/shape of the lossy volume; the bridge and port
dimensions and the port's position in the gap; the mesh and PML settings; the exact sensitivity and scan formalism
(the paper cites a DMRadio methods paper ``to appear'', 2022).

\paragraph{Reference data from the paper itself.} Lossless toroid: first self-resonance (pole of $Q_B$) at
26.7\,MHz; inductive-tuning bands 26.5--47.2, 75.1--84.0, 116.1--117.4\,MHz (extracted). Solenoid bands:
22.0--36.4, 69.7--92.5, 131.5--149.6\,MHz and several narrow ones.

\section{Status}
\begin{longtable}{p{2.2cm}p{4.6cm}p{6.6cm}p{2cm}}
\toprule figure & file & axionFEM inputs needed & status\\\midrule
%%TABLE%%
\bottomrule\end{longtable}

\clearpage
%%DIRECT%%

\section{Figures}
%%FIGURES%%

\end{document}
